%\pdfoutput=1

\documentclass[a4paper,12pt]{article}
%\usepackage{axodraw}
\usepackage{slashed}
%\usepackage[affil-sl,auth-sc]{authblk}
\usepackage{amssymb,amsmath,amsbsy}
\usepackage{mathrsfs}
%\usepackage{mathbbol}
\usepackage{bm}
\usepackage{epsfig}
%\usepackage{appendix}
\usepackage{hyperref}
\usepackage{mdframed}
\usepackage{adjustbox}
\usepackage[makeroom]{cancel}

%\usepackage{graphicx}
%\DeclareGraphicsRule{.tif}{png}{.png}{`convert #1 `dirname #1`/`basename #1 .tif`.png}
%\usepackage{cite}
%\usepackage{float}
\usepackage{xcolor}
 
%% This list environment is used for the references in the
%% Program Summary
%%
\newcounter{bla}
\newenvironment{refnummer}{%
\list{[\arabic{bla}]}%
{\usecounter{bla}%
 \setlength{\itemindent}{0pt}%
 \setlength{\topsep}{0pt}%
 \setlength{\itemsep}{0pt}%
 \setlength{\labelsep}{2pt}%
 \setlength{\listparindent}{0pt}%
 \settowidth{\labelwidth}{[9]}%
 \setlength{\leftmargin}{\labelwidth}%
 \addtolength{\leftmargin}{\labelsep}%
 \setlength{\rightmargin}{0pt}}}
 {\endlist}

\textheight=24cm
\textwidth=17cm
\topmargin=-2cm
\evensidemargin=-5mm
\oddsidemargin=\evensidemargin

%%%%%%%%%%%%%%%%%%%%%%%%%%%%%%%%
\def\code{{\tt SUSY\_FLAVOR}}
\def\MMI{{\tt MassToMI}}
\def\sfr{{\tt SmeftFR}~}
\def\frules{{\tt FeynRules}~}
\def\webpage{\url{www.fuw.edu.pl/smeft}}
%%%%%%%%%%%%%%%%%%%%%%%%%%%%%%%%%%

%%%%%%%%%%%%%%%%%%%%%%%%%%%%%%%
\setlength{\parskip}{1mm}
%%%%%%%%%%%%%%%%%%%%%%%%%%%%%%%

%%%%%%%%%%%%%% Citing equations %%%%%%%%%%%
\def\eq#1{eq.~(\ref{#1})}
\def\Eq#1{Eq.~(\ref{#1})}
\def\Eqs#1#2{Eqs.~(\ref{#1}) and (\ref{#2})}
\def\eqs#1#2{eqs.~(\ref{#1}) and (\ref{#2})}
\def\eqss#1#2#3{eqs.~(\ref{#1}), (\ref{#2}) and (\ref{#3})}
\def\eqst#1#2{eqs.~(\ref{#1})--(\ref{#2})}
\def\Eqst#1#2{Eqs.~(\ref{#1})--(\ref{#2})}
\def\Ref#1{ref.~\cite{#1}}
\def\Refs#1{refs.~\cite{#1}}

%%%%%%%%%%%%%%%%%%%%%%%%%%%%%%%%%%%%%
\newcommand{\bra}[1]{\langle #1|}
\newcommand{\ket}[1]{|#1\rangle}
\newcommand{\braket}[2]{\langle #1|#2\rangle}
\newcommand{\ie}{{\it i.e., }}
\newcommand{\eg}{{\it e.g., }}
\newcommand{\Lu}{{\cal L}}
\newcommand{\Ou}{{\cal O}}
\newcommand{\bea}{\begin{eqnarray}}
\newcommand{\eea}{\end{eqnarray}}
\newcommand{\nn}{\nonumber\\}
\newcommand{\MeV}{\; \mathrm{MeV}}
\newcommand{\GeV}{\; \mathrm{GeV}}
\newcommand{\TeV}{\; \mathrm{TeV}}

\newcommand{\bce}{\begin{center}}
\newcommand{\ece}{\end{center}}
\newcommand{\be}{\begin{equation}}
\newcommand{\ee }{\end{equation}}
\newcommand{\btb}{\begin{tabular}}
\newcommand{\etb}{\end{tabular}}
\newcommand{\f}{\frac}
\newcommand{\eps}{\varepsilon}
\newcommand{\vp}{\varphi}

\def\lcal{{\cal L}}
\def\ocal{{\cal O}}
\def\wt{\widetilde}

\newcommand{\tvp}{\widetilde{\varphi}}
\newcommand{\vpj }{\mbox{${\vp^\dag i\,\raisebox{2mm}{\boldmath
        ${}^\leftrightarrow$}\hspace{-4mm} D_\mu\,\vp}$}}
\newcommand{\vpjt}{\mbox{${\vp^\dag i\,\raisebox{2mm}{\boldmath
        ${}^\leftrightarrow$}\hspace{-4mm} D_\mu^{\,I}\,\vp}$}}
%


\newcommand{\osum}[1]{\vspace{0.5cm}{\pmb{\Circle}}{\hspace{-0.4cm}\sum_{#1}\hspace{0.1cm}} }
\newcommand{\oosum}[1]{\vspace{0.5cm}{\pmb{\Circle}}{\hspace{-0.49cm}\sum_{#1}\hspace{0.1cm}} }

\usepackage{color}
\newcommand{\tcl}{\textcolor}
\definecolor{orange}{rgb}{0.9,0.2,0}
\definecolor{brown}{rgb}{0.7,0.3,0.2}
\definecolor{fuxia}{rgb}{1,0,1}
\definecolor{skyblue}{rgb}{0,0.1,0.9}
\definecolor{violetred}{rgb}{0.8,0.13,0.56}
\definecolor{deeppink}{rgb}{1.00,0.08,0.5}
\definecolor{pink}{rgb}{1.00,0.75,0.80}
\definecolor{orchid}{rgb}{0.85,0.44,0.84}
\definecolor{lightpink}{rgb}{1.00,0.71,0.76}
\definecolor{bluish}{rgb}{0,0.6,0.8}
\definecolor{lightgray}{rgb}{0.95,0.95,0.95}

\newcommand{\red}[1]{\color{red} #1 \color{black}}
\newcommand{\blue}[1]{\color{blue} #1 \color{black}}


\numberwithin{equation}{section}


\begin{document}

\begin{center}
  {\Large \bf Adding fermionic dimension-8 operators to \sfr v4}
\end{center}

As dimension-8 fermionic operators are very numerous, we did not
include full set of them as a standard feature of published code, but
we provide instructions on how to do that, illustrating it below on a
chosen example of fermionic operator.  We are working in the basis
of~\cite{Murphy:2020rsh}.  Below we list the steps which one should
follow.
\begin{itemize}
  \item Define the Mathematica code for the operator itself.  We
    present this on the example from Table 7 in
    Ref.~\cite{Murphy:2020rsh}:
\begin{equation*}
Q_{e^2 W H^2 D}^{(4)}=\left(\bar{e}_p \gamma^\nu
e_r\right)\left(H^{\dagger} \overleftrightarrow{D}^{I \mu} H\right)
\widetilde{W}_{\mu \nu}^I
\end{equation*}

The corresponding definition should be placed in file {\tt
  Lagrangian/215\_psiXphi2D.fr} (we follow the numeration of operator
classes in~\cite{Murphy:2020rsh} - operator from class XX should be
defined in the file {\tt lagrangian/2XX\_class\_name.fr} ):
    
\bigskip
\noindent
    {\small \tt LQe2Wphi2Dn4 :=
      Module[\{sp1,sp2,m,ii,jj,ff1,ff2,mu,nu,al,be,aux\},

    aux = 2 Ta[m,ii,jj] (Phi8bar[ii] DC[Phi8[jj],mu] -
    DC[Phi8bar[ii],mu] Phi8[jj])*\\ lRbar[sp1,ff1].lR[sp2,ff2]
    Ga[nu,sp1,sp2] Eps[mu,nu,al,be]/2 HC[FS[Wi,al,be,m]];
    
    aux = ExpandIndices[ ToExpression[SMEFT\$WB <> "e2Wphi2Dn4"][ff1,ff2] aux,\\
        $\phantom{XXXXXXXXXXX}$
         FlavorExpand->\{SU2W,SU2D\} ];

    aux /. SMEFTGaugeRules ]; }

\bigskip
We follow standard FeynRules syntax (see SM Lagrangian definition
included as an example in FeynRules distribution), with the exception
of replacement of {\tt Phi $\to$ Phi8} used to speed up expansions in
mass dimensions.  For consistency and to avoid accidental name
repetitions, naming of operators in code should follow their naming
in~\cite{Murphy:2020rsh} (preceded by ``L'' as ``Lagrangian''), in
this case $Q_{e^2 W H^2 D}^{(4)}\to$ {\tt "e2Wphi2Dn4"} (in the
operator names we also replace {\tt H} $\to$ {\tt phi}), with
Lagrangian entry named {\tt LQe2Wphi2Dn4}, extra {\tt "LQ"} always
preceding operator name.

\item Open file {\tt/code/smeft\_variables.m} and:
\begin{enumerate}
\item Add operator name {\tt "e2Wphi2Dn4"} to the relevant list near
  the top of file, their names are self-explanatory:\\ {\tt
    SMEFT\$Dim8TwoFermionOperators,\\
    SMEFT\$Dim8FourLeptonOperators,\\
    SMEFT\$Dim8FourQuarkOperators,\\
    SMEFT\$Dim8TwoQuarkTwoLeptonOperators,\\
    SMEFT\$Dim8FourFermionBLVOperators}.\\[2mm]
  In our example obviously one should include it in\\
  {\tt SMEFT\$Dim8TwoFermionOperators=\{\dots,"e2Wphi2Dn4",\dots\} }.
\item Add new line to the {\tt SMEFT\$TensorWC} table:\\[2mm]
%
  {\tt SMEFT\$TensorWC = \{\\
%
  $\phantom{XXXXX}$  \dots\\
%
  $\phantom{XXXXX}$  \{"e2Wphi2Dn4",\{VLR,VLR\},False,2\},\\
    %
    $\phantom{XXXXX}$ \\
%
   $\phantom{XXXXX}$  \dots \}
  }
  \newpage
The syntax of the extra lines in {\tt SMEFT\$TensorWC} consist of:
    \begin{itemize}
    \item[$1^{st}$] entry: operator name
    \item[$2^{nd}$] entry: list of fermionic rotation matrices. They
      are denoted as {\tt VXY}, where {\tt X=(L,U,D)} is the fermion
      type and {\tt Y} its chirality. For instance, for operator
      containing fermionic current $l_L u_R$ one should add as a
      $2^{nd}$ entry {\tt \{VLL, VUR\}}. Similar entry for 4-fermion
      $q_L q_L d_R d_R $ operator would be {\tt \{VDL, VDL, VDR,
        VDR\}}, and so on. For left $SU(2)$ doublet $q_L$ either {\tt
        VUL} or {\tt VDL} can be used, it affects in which places CKM
      matrix will appear in fermion mass basis interactions.
    \item [$3^{rd}$] entry: is the operator $B$ or $L$ violating?
      Always {\tt False} for 2-fermion dimension-8 operators.
    \item [$4^{th}$] entry: operator ``symmetry class''. Here for
      details see comments on index symmetries of WCs of dimension-6
      fermionic operators (``Straub classes'') in file {\tt
        code/smeft\_variables.m}, above the definition of {\tt
        TensorInd} table. For dimension-8 fermionic operators, one can
      usually assign for them one of the classes defined for
      dimension-6 case. If flavor index symmetries for newly added
      operator are not known or differ from dimension-6 ones, one can
      assign class 1 for 2-fermion and class 9 for 4-fermion case,
      corresponding to no symmetries imposed (they can still be
      applied later at numerical level, initializing values of
      WCs). Experienced and brave user can also define new types of
      symmetry classes, modifying appropriately {\tt TensorInd} table
      and {\tt TensorWCSymmetrize} routine in {\tt
        code/smeft\_functions.m} file.
      \end{itemize}
\end{enumerate}
\item Go to {\tt smeft\_fr\_init.m} (or {\tt SmeftFR-init.nb}) and add
  operator name (here {\tt "e2Wphi2Dn4"}) to the {\tt OpList8} which
  lists dimension-8 operators included in calculations.
\end{itemize}


\begin{thebibliography}{1}
\expandafter\ifx\csname url\endcsname\relax
  \def\url#1{\texttt{#1}}\fi
\expandafter\ifx\csname urlprefix\endcsname\relax\def\urlprefix{URL }\fi
\expandafter\ifx\csname href\endcsname\relax
  \def\href#1#2{#2} \def\path#1{#1}\fi

\bibitem{Murphy:2020rsh} C.~W. Murphy, {Dimension-8 operators in the
  Standard Model Effective Field Theory}, JHEP 10 (2020) 174.  \newblock
  \href {http://arxiv.org/abs/2005.00059} {\path{arXiv:2005.00059}},
  \href {https://doi.org/10.1007/JHEP10(2020)174}
        {\path{doi:10.1007/JHEP10(2020)174}}.

\end{thebibliography}


\end{document}
